% Anhang / Zusatzmaterial: zählt nicht zu den 6 Seiten (Folie 16). Format wie im Author-Kit
% (sec/X_suppl.tex), Überschrift auf Deutsch. Zahlen aus results/ms1, results/why, results/explain.
\clearpage
\setcounter{page}{1}
\twocolumn[\centering\Large\textbf{\thetitle}\\\vspace{0.5em}Anhang\\\vspace{1.0em}]
\appendix

\section{Hyperparameter}
\label{sec:hyperparameters}
\Cref{tab:hyper} listet alle Konfigurationen (\texttt{configs/ms1/*.yaml}). Gemeinsam für alle: AdamW,
Weight Decay 0,05, Cosine-Schedule mit 10\,\% Warm-up, bis zu 10 Epochen, Early Stopping (Geduld 3),
Batchgröße 32, $224\times224$ Pixel Eingabe, Augmentierung \emph{light}, Seeds 1--3.

\begin{table}[h]
  \centering
  \caption{Lernrate und trainierbare Parameter je Methode.}
  \label{tab:hyper}
  \small
  \resizebox{\linewidth}{!}{\begin{tabular}{lll}
    \toprule
    Methode & Lernrate & Trainierbar \\
    \midrule
    ViT Linear Probing & $10^{-3}$ & Kopf \\
    ViT LoRA $r\in\{4,8,16\}$, $\alpha=2r$, Dropout 0,1, Query/Value & $5\cdot10^{-4}$ & LoRA + Kopf \\
    ViT LoRA $r=8$, alle linearen Layer & $5\cdot10^{-4}$ & LoRA + Kopf \\
    ViT volles Fine-Tuning & $5\cdot10^{-5}$ & alle \\
    ResNet-50, ConvNeXt-T & $2\cdot10^{-4}$ (Validation) & alle \\
    \bottomrule
  \end{tabular}}
\end{table}

\section{Warum sind die Accuracies so hoch?}
\label{sec:why}
Alle Zahlen aus \texttt{results/why/REPORT.md} (\texttt{make why}).

\paragraph{Pretraining.} Zufällig initialisiert und mit denselben Daten und derselben Auswertung
trainiert (30 Epochen, größere Lernrate), erreicht ViT-B/16 $11{,}18\pm2{,}61$\,\% und ResNet-50
$31{,}07\pm0{,}41$\,\% Test-Accuracy, gegenüber $95{,}39\pm0{,}39$ und $94{,}22\pm0{,}85$\,\% mit
ImageNet-Gewichten (Zufall: 4\,\%).

\paragraph{Vorwissen.} 22 der 25 Rassen sind ImageNet-1k-Klassen~\cite{russakovsky2015imagenet}; es
fehlen nur American Bulldog, Havanese und Shiba Inu. Ganz ohne Training klassifizieren die
ursprünglichen ImageNet-Köpfe unserer Checkpoints die Testbilder dieser 22 Rassen (Auswahl unter den 22)
mit 96,20\,\% (ViT-B/16), 96,75\,\% (ResNet-50) und 96,52\,\% (ConvNeXt-T). LoRA $r=4$ erreicht auf
denselben Bildern 95,88\,\%, bei Auswahl unter allen 25 Rassen. Die Zuordnung von Rasse zu
ImageNet-Klasse erfolgt über den Namen (Tabelle in \texttt{animal\_vit/why.py}); sie ist eine Annahme. Der
American Pit Bull Terrier wird der Klasse ``American Staffordshire terrier'' zugeordnet, Japanese Chin
``Japanese spaniel'', Scottish Terrier ``Scotch terrier'' und German Shorthaired ``German short-haired
pointer''. Bei diesen Rassen kann die Zuordnung strittig sein.

\paragraph{Eingefrorene Merkmale.} Eine logistische Regression~\cite{pedregosa2011sklearn} auf eingefrorenen Merkmalen (alle
Trainingsbilder) erreicht 95,94\,\% (ViT-B/16) und 96,06\,\% (ResNet-50), gegenüber 95,39 und 94,22\,\%
nach vollem Fine-Tuning. Mit einem Bild pro Rasse erreicht sie $84{,}66\pm2{,}63$ und
$87{,}41\pm1{,}91$\,\% (\cref{fig:fewshot} im Hauptteil).

\paragraph{Fast-Duplikate.} Für jedes Testbild haben wir die Kosinus-Ähnlichkeit zum ähnlichsten Trainings-
oder Validierungsbild mit Embeddings von CLIP ViT-B/16~\cite{radford2021clip} berechnet. 215 Testbilder
(8,6\,\%) haben einen Nachbarn mit Ähnlichkeit $\geq0{,}95$; ohne sie ändert sich die Accuracy von LoRA
$r=4$ von 96,11 auf 95,77\,\%.

\paragraph{Rest-Fehler.} Die vier Bully-Rassen (American Bulldog, American Pit Bull Terrier,
Staffordshire Bull Terrier, Boxer) machen 16\,\% der Testbilder und 59\,\% aller Fehler von LoRA $r=4$
aus (Accuracy 84,9\,\% gegenüber 98,2\,\% bei den übrigen Rassen). 71 Testbilder sind in allen drei
Seeds falsch klassifiziert.


\section{Warum diese Rasse und nicht die andere?}
\label{sec:explain}
Erklärungen für LoRA $r=4$, ConvNeXt-T und ResNet-50 (Seed 1; \texttt{results/explain/REPORT.md},
\texttt{make explain}). Wir nutzen Grad-CAM~\cite{selvaraju2017gradcam} (Implementierung:
pytorch-grad-cam~\cite{gildenblat2021gradcam}), Occlusion~\cite{zeiler2014visualizing} (Implementierung:
Captum~\cite{kokhlikyan2020captum}) und Attention Rollout~\cite{abnar2020rollout} (eigener Code). Ein
\emph{kontrastives} Grad-CAM verwendet das Ziel $\text{logit}_A-\text{logit}_B$ und hebt Evidenz für
Rasse $A$ und gegen Rasse $B$ hervor (vgl.~\cite{prabhushankar2020contrastive}). In einem
Deletion-Test~\cite{petsiuk2018rise} werden die auffälligsten $16\times16$-Patches zuerst grau
übermalt; Occlusion senkt die Klassenwahrscheinlichkeit beim ViT und bei ConvNeXt-T am schnellsten
(Fläche unter der Kurve 0,64 und 0,59 gegenüber 0,92 und 0,89 beim zufälligen Löschen). Gemessen mit der
jeweils treuesten Methode liegt die Saliency in der annotierten Kopf-Box (Validierungssplit) bei allen
drei Modellen beim 3,1- bis 3,7-Fachen ihres Flächenanteils. Grad-CAM liefert für den ViT einen deutlich
kleineren Wert (1,3), passend zu bekannten Grenzen von Grad-CAM bei Transformern~\cite{chefer2021transformer}.
\Cref{fig:head} im Hauptteil zeigt Beispielkarten mit den Kopf-Boxen, \cref{fig:contrastive} kontrastive Karten für
die häufigsten Verwechslungen.




\begin{figure}[t]
  \centering
  \resultfigure{explain/contrastive_convnext_t.png}{\linewidth}{Kontrastive Karten}
  \caption{Warum Rasse $A$ und nicht $B$? Für die häufigsten Verwechslungen: Grad-CAM für $A$ und $B$,
  kontrastive Karte $A-B$ und Occlusion für $A$ (ConvNeXt-T, Testbilder).}
  \label{fig:contrastive}
\end{figure}

\section{Software und KI-Hilfsmittel}
\label{sec:tools}
Die Experimente nutzen PyTorch~\cite{paszke2019pytorch}, Hugging Face
Transformers~\cite{wolf2020transformers}, scikit-learn~\cite{pedregosa2011sklearn}, pytorch-grad-cam~\cite{gildenblat2021gradcam}
und Captum~\cite{kokhlikyan2020captum}.
% Erklärung zur KI-Nutzung (Modulregeln, Vorlesung 01, Folie 13); in der Präsentation mit dem Dozenten besprochen.
Bei Code, Auswertung und Formulierung haben wir einen KI-Assistenten eingesetzt, wie in der Präsentation mit
dem Dozenten besprochen. Forschungsfragen, Experimente und die Prüfung aller Ergebnisse stammen von uns;
alle Zahlen stammen aus den protokollierten Experimenten (\texttt{results/}).
